\section*{ID6855: Radial Coordinate System}
\paragraph{Metadata.}
PiID: 8670543154706 | RTEID: RTE:P26619.4873:T6.1616 | UUID: 9e6291c5-fb4f-5a3b-bbdf-8e06f85b0bf6

\paragraph{Canonical Statement.}
\textbackslash{}[ R=\textbackslash{}sqrt\{X\textasciicircum{}2+Y\textasciicircum{}2\} \textbackslash{}]

\paragraph{Canonical Equation.}
\[
R=\sqrt{X^2+Y^2}
\]

\paragraph{Dictionary.}
Radial Coordinate System — R=√(X\textasciicircum{}2+Y\textasciicircum{}2)

\paragraph{Encyclopedia.}
Radial Coordinate System is an algebraic definition, written R=√(X\textasciicircum{}2+Y\textasciicircum{}2). It is an algebraic definition expressing the quantity directly in terms of the variables on its right-hand side. It is built from X, Y, defining R. Within the Trawin Zero Point registry it serves as one element of the framework's mathematical narrative.
